\section{Equal-rate descent and affine criticality}
This module preserves the geometric propositions of the affine-criticality note~\cite{SA}. They do not depend on the BFGS principal-minor lemma. The originally proposed dynamical step is now proved only in the completed settings of Modules C and U; no blanket theorem for all finite smooth maxima is asserted.

For a finite convex $F$, put
\[
 V(p)=\operatorname{span}(\partial F(p)-\partial F(p)),\quad U(p)=V(p)^\perp,
 \quad a(p)=\operatorname{proj}_{\operatorname{aff}\partial F(p)}0.
\]
The affine hull equals $a(p)+V(p)$ and $a(p)\perp V(p)$. Hence every subgradient satisfies
\[
 P_{U(p)}g=a(p),\qquad g^Ta(p)=\|a(p)\|^2.
\]
The condition
\begin{equation}\label{A:AC}
 0\in\operatorname{aff}\partial F(p)\Longrightarrow0\in\partial F(p)
\end{equation}
is therefore equivalent, at every nonoptimal point under consideration, to the existence of a direction $v$ with $g^Tv=-1$ for every $g\in\partial F(p)$. One may take $v=-a(p)/\|a(p)\|^2$. Conversely such an equal-rate direction rules out zero in the affine hull. This is stronger than ordinary strict descent, which requires only negative, not equal, inner products.

On a compact admissible run, all cluster points share the limiting value $\ell$. If at least one cluster point is affine critical and \eqref{A:AC} holds on the basin, that point is optimal, so $\ell=F_*$. Compactness gives convergence to the solution set, and uniqueness gives convergence of the entire sequence. This reduction does not assume an iterate limit. It would be incorrect to replace affine criticality of a cluster point by vanishing gradients at differentiable iterates: for example, the gradients of $|x|$ have unit magnitude away from its minimizer.

\section{Why pointwise affine separation need not be uniform}
Let $\gamma(p)=\operatorname{dist}(0,\operatorname{aff}\partial F(p))$. A useful quantitative condition is a positive minimum of $\gamma$ on every compact positive objective band. Pointwise nonvanishing plus compactness does not establish it when ranks change.

Consider
\[
 F(u,v)=u+u^2+v^2+|uv|
 =\max\{u+u^2+v^2+uv,\ u+u^2+v^2-uv\}.
\]
Both branch Hessians have eigenvalues $1$ and $3$, so $F$ is strongly convex, with unique minimizer $(-1/2,0)$ and value $-1/4$. At $(0,t)$, $t\ne0$, the active gradients are $(1+t,2t)$ and $(1-t,2t)$, giving $\gamma(0,t)=2|t|$. At zero the gradient is $(1,0)$, so $\gamma(0,0)=1$. On the other switching axis $\gamma(u,0)=|1+2u|$, vanishing only at the optimizer. Smooth points are affine critical only when stationary. Thus \eqref{A:AC} holds globally, but $\gamma(0,t)\to0$ while the objective gap tends to $1/4$.

This is a geometric rank-loss example, not a BFGS counterexample. A sufficiently small neighborhood of its actual optimizer excludes the rank-loss point.

\section{Nondegenerate smooth maxima give local affine-critical correctness}
Let $F=\max_{i\le m}f_i$ with convex $C^1$ branches and a minimizer $x_*$. Restrict to $I_*$, the active branches there; inactive branches have positive gaps and disappear in a sufficiently small neighborhood. Suppose the active gradients are affinely independent and admit strictly positive coefficients
\[
 \lambda_i^*>0,\quad\sum_{i\in I_*}\lambda_i^*=1,
 \quad\sum_{i\in I_*}\lambda_i^*\nabla f_i(x_*)=0.
\]
Then \eqref{A:AC} holds near $x_*$.

\begin{proof}
Let $A(x)$ have columns $(\nabla f_i(x),1)$ for $i\in I_*$. Full column rank at $x_*$ persists nearby. At an affine-critical point $p$, extend the affine coefficients of its active gradients by zero on missing indices. They solve
\[
 A(p)\lambda=(0,\ldots,0,1)^T.
\]
Whenever this system is solvable, its unique solution equals
\[
 \lambda(p)=(A(p)^TA(p))^{-1}A(p)^T(0,\ldots,0,1)^T.
\]
This continuous expression equals the strictly positive vector $\lambda^*$ at $x_*$. It remains positive nearby. No missing index can therefore have zero coefficient, and the affine combination is a convex combination. Thus $0\in\partial F(p)$.
\end{proof}
This is the source note's barycentric-stability argument, closely related to the geometric nondegeneracy results it cites. It is not itself a BFGS convergence theorem.

There is also local lower semicontinuity of $\gamma$. For each nonempty $J\subset I_*$, the distance to the affine hull of the gradients in $J$ is continuous, because the relevant difference Gram matrix remains nonsingular. For $p_j\to p$, pass to a subsequence of constant active set $J\subset I(p)$. Then
\[
 \lim\gamma(p_j)=\operatorname{dist}(0,\operatorname{aff}\{\nabla f_i(p):i\in J\})\ge\gamma(p).
\]
On a compact positive band, local \eqref{A:AC} therefore gives a strictly positive minimum. This establishes visibility, not the full transport needed in arbitrary dimension.

\section{Semialgebraic objectives have finitely many affine-critical values}
\begin{proposition}[A low-level affine-critical gap]\label{A:semialgebraic}
If $F:\mathbb R^n\to\mathbb R$ is finite convex and semialgebraic and attains $F_*$, then some $\epsilon>0$ satisfies
\[
 F(p)<F_*+\epsilon,\quad0\in\operatorname{aff}\partial F(p)
 \quad\Longrightarrow\quad0\in\partial F(p).
\]
\end{proposition}
\begin{proof}
The subgradient graph is semialgebraic by its quantified supporting-plane description and quantifier elimination. Using at most $n+1$ subgradients in an affine representation shows that the affine-critical set $C_{\rm aff}$ is semialgebraic. It has finitely many semialgebraically path-connected components.

On a piecewise $C^1$ semialgebraic path $p(t)$ in that set, the composition $F\circ p$ is continuous and piecewise $C^1$. At a time where both derivatives exist, the two-sided convex subgradient inequality implies
\[
 (F\circ p)'(t)=g^Tp'(t)\qquad\text{for every }g\in\partial F(p(t)).
\]
Taking an affine combination of these subgradients equal to zero proves $(F\circ p)'(t)=0$. Thus $F$ is constant on each path component. There are finitely many affine-critical values. Since $F_*$ is one of them, choose $\epsilon$ below the least strictly larger value, or any positive number if no larger value occurs. Convexity turns equality with $F_*$ into stationarity.
\end{proof}
The standard semialgebraic path and projection facts are the external geometric ingredients cited in~\cite{SA}, including Bolte--Daniilidis--Lewis--Shiota~\cite{BDLS}. This proposition concerns affine criticality rather than merely Clarke criticality. It does not imply a blanket local BFGS theorem without a dynamical result forcing such a cluster point. Modules C and U provide completed restricted cases.

\section{What remains a conjecture}
The recorded general-dimensional conjecture concerns finite convex $C^{2,1}$ maxima with unique minimizer, compact low sublevels, affinely independent active gradients, positive optimal weights, and positive weighted Hessian on
\[
 U_* =\{v:(\nabla f_i(x_*)-\nabla f_1(x_*))^Tv=0\ \forall i\in I_*\}.
\]
The proposed conclusion is local full-BFGS convergence for arbitrary SPD initialization. The planar theorem proves a strict-contour version; the fixed-normal theorem proves a different arbitrary-dimensional class. Neither establishes this full conjecture. Straightening an active manifold by a nonlinear coordinate change does not preserve ordinary BFGS, and cannot supply the missing implication.
